\section{The two frame--time couplings}
\label{sec:couplings}

Theorem~\ref{thm:classify} tells us \emph{which} directions are undetermined; this section
quantifies \emph{how strongly} the frame and the timescale couple along two physically
distinct channels. The classical roots of both channels are \emph{cited}, not claimed; the
theorem structure around them is \Proven{}; and the magnitudes are \Modelled{}
(first-principles, cross-checked, not certified) and tagged \modtag{}. One of the two
subsections corrects a claim whose naive form is false, and we present the correction
openly as the honest result.

\subsection{Coupling 1: range \texorpdfstring{$\leftrightarrow$}{<->} clock offset}
\label{sec:coupling1}

\paragraph{The unit tie and the classical root.}
For one-way ranging, a radial position shift $\delta d$ and a clock offset $c\,\dtau$ enter
every measurement with the same sign: a range is lengthened equally by moving the target
away or by advancing the clock. The exact exchange rate is $c$,
\begin{equation}
  1\,\text{m}\;\longleftrightarrow\;\frac{1}{c}\;=\;3.336\,\text{ns},
  \label{eq:unittie}
\end{equation}
computed from the engine's $c$ (\code{lunar\_time::C2\_M2\_S2}). This is precisely the
GNSS radial-position/receiver-clock degeneracy, classical in the integrity and positioning
literature \citep{sturza1988,brown1992,parkinson1988}; we cite it and claim no novelty. Its
importance here is that it is the strongest frame--time coupling, and the natural place to
look for a \emph{coupled datum defect}.

\paragraph{It is a near-defect, not an exact one.}
Crucially, \eqref{eq:unittie} is a \emph{near}-defect: for a single user with poor local
vertical geometry the scale/radial and offset signatures are nearly parallel, giving a
small---but strictly positive---eigenvalue, not a member of $\ker\Imat$. It becomes an
exact null direction only in the degenerate limit of perfectly co-directional geometry.
Accordingly it does \emph{not} contribute to the defect $d$ of Theorem~\ref{thm:classify};
we quantify it instead by the $2\times2$ Schur-complement (marginal) Fisher onto
$\{s=\text{idx }3,\ \dtau=\text{idx }7\}$, eliminating the other seven parameters
(\code{coupled\_marginal\_fisher} / \code{coupled\_marginal\_eigs}; the algebra is in
Appendix~\ref{app:methods}). The smallest eigenvalue $\lambda_{\min}$ of this marginal is
the coupling-1 near-defect strength: the degree to which the lunar body's scale
(radial breathing) is observable \emph{independently} of the common clock offset in a given
one-way network.

\paragraph{The corrected, defensible finding.}
The tempting design claim is: ``two-way ranging (which cancels the node clock, cf.\
\eqref{eq:twoway}) selectively breaks the scale$\leftrightarrow$offset trade and lifts the
near-defect, so a lunar system needs two-way ranging to separate frame scale from
timescale offset.'' On a real-code-path probe this claim proves \emph{false at the datum
level}, and we report the correction. On a geometrically diverse multi-beacon
multi-epoch one-way network (five beacons $\times$ six orbiters $\times$ four epochs,
$120$ one-way rows, real DE440 orientation), the coupled Fisher is \emph{full-rank}
($d=0$) and the $\{s,\dtau\}$ marginal is positive-definite: its eigenvalues are
$[\,1.872,\ 36.00\,]$ (ascending), both strictly positive. Scale is already observable
independently of the offset. We then compare, at equal observation count, two ways of
adding $120$ rows to this base (Table~\ref{tab:coupling1}): (a) $120$ two-way rows at the
same geometry, and (b) $120$ \emph{more} one-way rows at a jittered geometry (argument of
latitude shifted by $7^\circ$). Both raise $\lambda_{\min}$ of the marginal; the ratio of
the two lifts is
\begin{equation}
  \frac{\text{lift}_{\text{two-way}}}{\text{lift}_{\text{more one-way}}}\;\approx\;1.00,
  \label{eq:liftratio}
\end{equation}
altitude-independent within the tested networks. In words: \textbf{in a geometrically
diverse network, two-way ranging provides no selective scale$\leftrightarrow$offset
decoupling beyond data volume.} The strong $1\,\text{m}\leftrightarrow3.336\,\text{ns}$
degeneracy is a \emph{per-node}, single-user, poor-local-geometry effect (the classical
GNSS radial/clock degeneracy) that does \emph{not} aggregate into a datum-level near-defect
once the beacon and orbiter directions are geometrically diverse.

\begin{table}[t]
\centering
\caption{Coupling~1 at the datum level (\Modelled{} geometry, real DE440 orientation;
\modtag{}). On a geometrically diverse $120$-row one-way network the coupled Fisher is
full-rank and the $\{s,\dtau\}$ Schur marginal is positive-definite. Adding $120$ two-way
rows raises the marginal $\lambda_{\min}$ by the \emph{same} factor as adding $120$ more
one-way rows: two-way ranging gives no selective decoupling beyond data volume. The
marginal eigenvalue pair $[1.872,\,36.00]$ is reproduced by the independent numpy/LAPACK
oracle (Section~\ref{sec:anchor}).}
\label{tab:coupling1}
\small
\begin{tabular}{@{}lcc@{}}
\toprule
\textbf{Network} & $\{s,\dtau\}$ marginal $\lambda_{\min}$ & \textbf{status} \\
\midrule
Base: $120$ one-way rows (diverse, defect $0$) & $1.872$ & \Validated{} (Table~\ref{tab:anchor}) \\
$+\,120$ two-way rows (same geometry)          & lift $>1$ & \Modelled{} \\
$+\,120$ more one-way rows (jittered $7^\circ$) & lift $>1$ & \Modelled{} \\
lift ratio (two-way / more one-way)             & $\approx1.00$ & \Modelled{} \emph{(the finding)} \\
\bottomrule
\end{tabular}
\end{table}

\begin{remark}[The idealized single-node regime, for contrast]
\label{rem:singlenode}
The naive claim is not \emph{empty}: in an idealized single-node, near-rank-one regime
(one beacon, poor vertical geometry) a two-way tie \emph{does} raise $\lambda_{\min}$ by
orders of magnitude, because there the scale and offset signatures really are nearly
parallel. That regime is a unit-test fixture, not a lunar constellation. The design-relevant
statement is the diverse-network one \eqref{eq:liftratio}; we keep the two regimes strictly
separate, and neither is a coupled \emph{defect} (Remark~\ref{rem:distinct}).
\end{remark}

This yields the first design law, developed in Section~\ref{sec:designlaws}:
\emph{geometric diversity separates the lunar frame scale from the timescale offset in
one-way ranging; a two-way tie or an external time tie is required only in the
low-diversity / single-user regime.}

\subsection{Coupling 2: relativistic rate \texorpdfstring{$\leftrightarrow$}{<->} frame}
\label{sec:coupling2}

The second channel is the dependence of a lunar-surface clock's \emph{rate} on the
frame it is realized in. A clock runs at a rate set by the gravitational potential and the
squared speed of its location (relativistic geodesy and the IAU relativistic timescale
framework \citep{soffel2003iau,mueller2018clocks}); perturbing the frame datum---its scale,
a radial position, or an inferred velocity---perturbs that rate. This is a real coupling,
and we ask only how \emph{strong} it is.

\paragraph{The frame\,$\to$\,rate Jacobian (\Modelled{}, \modtag{}).}
Writing $\Umoon=GM_{\mathrm{M}}/\rmoon$ for the lunar surface potential and
$\gmoon=GM_{\mathrm{M}}/\rmoon^2$ for the surface gravity, the engine's
\code{rate\_frame\_jacobian} returns the three first-principles coefficients
\begin{align}
  \frac{\partial\alpha}{\partial s}
    &\;=\; +\frac{\Umoon}{c^2}
    \;=\; +3.140\times10^{-11}\ \text{per unit fractional scale},
    \label{eq:dads}\\
  \frac{\partial\alpha}{\partial v}
    &\;=\; -\frac{|v|}{c^2}
    \;\approx\; -1.11\times10^{-14}\ \text{per (m/s)},
    \label{eq:dadv}\\
  \frac{\partial\alpha}{\partial r_{\mathrm{radial}}}
    &\;=\; +\frac{\gmoon}{c^2}
    \;=\; +1.807\times10^{-17}\ \text{per metre}.
    \label{eq:dadr}
\end{align}
The radial and scale coefficients satisfy the exact identity
$\partial\alpha/\partial r=(\partial\alpha/\partial s)/\rmoon$ (unit-tested to
$10^{-12}$), since $\gmoon/c^2=(\Umoon/\rmoon)/c^2$. The velocity coefficient is
epoch-dependent because the geocentric Moon speed varies by several percent over the lunar
month (the $969$ and $997$~m/s values below already differ by nearly $3\%$); the engine
prints $-1.078\times10^{-14}$ at the fixture epoch 2024-01-01
($|v|\approx969$~m/s) versus $\approx-1.11\times10^{-14}$ near J2000
($|v|\approx997$~m/s). We quote both as honest epoch variation, not an error. The signs
combine as
\begin{equation}
  \dalpha
  \;=\; \frac{\partial\alpha}{\partial s}\,\delta s
  \;+\; \frac{\partial\alpha}{\partial r}\,\delta r
  \;-\; \Bigl|\frac{\partial\alpha}{\partial v}\Bigr|\,|\delta v|,
  \label{eq:dalphacombine}
\end{equation}
the velocity term subtracting (a faster clock runs slower). The velocity and radial
couplings are physical but, as noted in Section~\ref{sec:rows}, do not map onto the
nine-vector datum; only the scale coupling appears as the datum-level rate tie
\eqref{eq:ratetie}.

\paragraph{Propagating an ILRF-class frame datum to the rate (\Modelled{}, \modtag{}).}
Take frame-tie uncertainties at the level a lunar analogue of the terrestrial reference
frame realization would achieve: a scale tie $\delta s\approx6.9\times10^{-8}$
(i.e.\ $0.12$~m of radial equivalent, $0.12/\rmoon$), a velocity of
$\delta v\lesssim1$~mm/s, and a radial position of $\delta r\lesssim0.1$~m. Through
\eqref{eq:dads}--\eqref{eq:dadr} these propagate to
\begin{equation}
  \dalpha(\text{scale})\approx2.17\times10^{-18},\quad
  \dalpha(v,\,1\,\text{mm/s})\approx1.11\times10^{-17},\quad
  \dalpha(\text{radial},\,0.1\,\text{m})\approx1.81\times10^{-18},
  \label{eq:propagate}
\end{equation}
so the frame \emph{realization} perturbs the lunar timescale rate by at most a
frame-realization envelope
\begin{equation}
  \boxed{\;\dalpha_{\text{frame realization}}\;\lesssim\;\sim2\times10^{-17}\;}
  \qquad\text{(realization envelope, kinetic-dominated).}
  \label{eq:framebound}
\end{equation}
These two budgets must be kept apart. Only the scale term of \eqref{eq:propagate} is a
genuine \emph{datum}-channel gain: $\dalpha(\text{scale})\approx2.2\times10^{-18}$ is what a
nine-vector datum-defect $\delta s$ propagates to the rate (Section~\ref{sec:designlaws}).
The larger velocity and radial contributions are frame-\emph{realization} /
orbit-determination errors that, as \eqref{eq:ratetie} and Section~\ref{sec:rows} note, do
\emph{not} map onto the datum; the envelope \eqref{eq:framebound} merely absorbs them for a
conservative bound. The gauge-robustness conclusion holds \emph{a fortiori}---the true
datum-channel gain is an order of magnitude below the realization envelope. Nothing larger
exists in either budget: the tidal-gradient contribution is $\sim2.7\times10^{-22}$ per
metre, and the external/tidal potential is excluded from $\Umoon$ by the equivalence
principle (a freely-falling body does not sense the uniform external potential, so we do not
add $GM_{\oplus}/d$). We quote the envelope as $\lesssim\!\sim\!2\times10^{-17}$, not
$\lesssim10^{-17}$, to keep the kinetic term honest.

\paragraph{The reference-surface definition band (\Modelled{}, \modtag{}).}
Against this estimation coupling stands the \emph{definitional} freedom of the lunar
timescale rate---the choice of reference surfaces, reference potential $W_0$, and averaging
window that a standard must fix. The lunar-clock-rate band relative to TT spans
$[56,59]$~$\mu$s/day, consistent with the lunar-timescale figures discussed in the Moonlight
reference-system literature
% confirm the [56,59] $\mu$s/day figure against fienga2024moonlight at camera-ready
\citep{fienga2024moonlight}; its half-width $\pm1.5$~$\mu$s/day
converts to a dimensionless rate
\begin{equation}
  \pm\frac{1.5\times10^{-6}\ \text{s/day}}{86400\ \text{s/day}}
  \;=\;\pm1.736\times10^{-11},
  \label{eq:band}
\end{equation}
about $10^{6}$ times the estimation coupling \eqref{eq:framebound}. The two axes are
orthogonal: one is set by \emph{how well the frame is estimated}, the other by \emph{how the
reference surfaces are defined}. The band's physical content includes named contributors
inside the lunar potential---tidal deformation $\approx2.4\times10^{-16}$ and $J_2$
oblateness $\approx6\times10^{-15}$, together with the periodic solar tidal contribution to
the Earth--Moon system $\approx6\times10^{-14}$ (periodic, and \emph{larger} than the $J_2$
term)---this named list is illustrative, not exhaustive (dimensionless; physics-literature
reference values, not engine-computed); its full $\pm1.7\times10^{-11}$ width is dominated by
the definitional $W_0$/selenoid-surface choice and the averaging window. A model-center
cross-check reconciles the band midpoint from first principles: the Earth-geoid redshift
that defines TT contributes $\approx60.2$~$\mu$s/day, the lunar surface potential subtracts
$\approx2.71$~$\mu$s/day, and the lunar orbital kinetic term subtracts
$\approx0.48$~$\mu$s/day, giving $60.2-2.71-0.48=57.0$~$\mu$s/day, comfortably within
$[56,59]$ (all \Modelled{}, first-principles).

\paragraph{Where the estimation coupling would matter.}
The estimation bound \eqref{eq:framebound} of $\sim2\times10^{-17}$ sits below the
stability of near-term deployed optical clocks ($\sim10^{-16}$) and would become relevant
only for laboratory-grade clocks at $\sim10^{-18}$. It is, for any foreseeable operational
lunar timescale, negligible against the definitional band.

\paragraph{Headline (\Proven{} logic, \Modelled{} magnitudes).}
The lunar timescale rate is \textbf{gauge-robust to the frame estimation datum}
($\lesssim\!\sim\!2\times10^{-17}$) but \textbf{gauge-sensitive to the reference-surface
definition} ($\pm1.7\times10^{-11}$). The dominant gauge of the rate is
\emph{definitional}, not \emph{estimational}: no PNT observation reduces it, and a standard
must \emph{fix} the surfaces and conventions. This is the time-analogue of the spatial
result of \citep{baweja2026datum}---the datum is fixed by convention and external ties, not
estimated---transported to the rate of the timescale.
